\section{Conclusion}
This study compared data storytelling enhanced visualisations to conventional visualisations across tasks of varying complexity based on Bloom’s taxonomy. In particular, we focused on line charts and choropleth maps to communicate the data stories and support task completion. The results show that the effectiveness and efficiency are not significantly affected in the simpler tasks like extracting values from the graphs (\textit{Identify}). Whereas, it does improve the effectiveness in \textit{Understanding} tasks and the efficiency in higher-thinking tasks ranging from \textit{Understanding} to \textit{Evaluating}. Interestingly, DS-enhanced visualisations reduced the number of \textit{Evaluate} sentences in creation tasks, suggesting that they may decrease critical thinking by making users overly confident in the presented data. This highlights potential risks when using DS-enhanced visualisations in sensitive contexts, such as news or policy discussions, where critical evaluation is crucial. 
Overall, most participants preferred DS-enhanced visualisations for their clarity and contextual information, while some favoured the simplicity of conventional visualisations. These findings suggest that visualisation design should be adaptable to user preferences, task complexity, and task goals.

